\chapter{Praktischer Kontext und Anforderungen}
\label{chap:praxiskontext-anforderungen}

\section{Betrieblicher Hintergrund}
\label{sec:betrieblicher-hintergrund}
In betrieblichen IT-Landschaften nehmen Webservices und Microservices eine zentrale Rolle ein. Ein typisches Szenario stellt die automatisierte Auftrags- und Bestellverarbeitung (Order Processing API) dar, bei der Clients Aufträge einreichen, Statusabfragen durchführen und asynchrone Folgeprozesse angestoßen werden. Dienstunterbrechungen führen in solchen Systemen unmittelbar zu wirtschaftlichen Einbußen und Reputationsverlust.

\section{Anforderungen an die Betriebsumgebung}
\label{sec:anforderungen-betriebsumgebung}
Aus dem betrieblichen Szenario leiten sich spezifische funktionale und nicht-funktionale Anforderungen an die Plattform ab:
\begin{itemize}
    \item \textbf{Reproduzierbarkeit:} Die gesamte Infrastruktur muss als Code versioniert und mit minimalem manuellem Eingriff reproduzierbar aufgebaut und zerstört werden können.
    \item \textbf{Transparenz des Systemzustands:} Aktuelle Versionen, Pod-Identitäten und Antwortzeiten müssen für Monitoring und Fehlersuche unmittelbar ersichtlich sein.
    \item \textbf{Fehlertoleranz:} Einzelne Instanzausfälle dürfen nicht zur Nicht-Erreichbarkeit des Gesamtsystems führen.
    \item \textbf{Kosteneffizienz:} Der Betrieb muss innerhalb eines strikten Budgets (unter 30~EUR pro Monat) realisierbar sein.
\end{itemize}

\section{Anforderungen an den Praxistransfer}
\label{sec:anforderungen-praxistransfer}
Der Praxistransfer erfordert, dass die im Rahmen dieser Arbeit gewonnenen Erkenntnisse und Implementierungsmuster generalisierbar und direkt auf vergleichbare betriebliche Anwendungen übertragbar sind. Die gewählten Mechanismen (Health Checks, GitOps-Sync, Terraform-Module) dürfen keine proprietären Abhängigkeiten erzeugen, die einen Transfer in reguläre Unternehmensumgebungen verhindern.
